\section*{Abstract}
Raylo's live Open Banking risk model reads Plaid's category names directly, which leaves a core model input outside our control. This report describes the taxonomy built to replace that dependency (275 detailed categories rolling up to 29 general ones), the seven-tier pipeline that applies it, and the modelling work completed since the August report.

On 2,000 held-out transactions the full pipeline now assigns the correct detailed category 83\% of the time. The provider's own category, which is what the live model consumes, is right 30\% of the time on the same rows. For the machine-learning tier that handles the transactions rules cannot, a transformer pretrained on 21.5 million bank narratives and taught from 405,000 consensus-labelled texts beats the linear classifier by five points on never-seen merchants and by ten points on high-risk categories, with confidence intervals that exclude zero.

On the credit side, a champion model built on leaf-level taxonomy features reaches a GINI of 0.62 on the six-month arrears outcome against 0.42 for the live model, up from the 0.56 quoted in August. Most of that gain comes from the model seeing thousands of leaf-level columns. Held to a governable 50 features and a single XGBoost, the same recipe scores 0.59, and that capped model is the one carried forward. A separate programme tested whether a sequence transformer reading raw transactions could beat it; on its own it could not. Added to the 50-feature model as single columns, however, a text-encoder score and a sequence-encoder score each give a clear lift, and together take the model to 0.63, above the uncapped champion. The text-score recipe is locked; the two-score version is the registered challenger for the prospective test.

